\section*{Chapter \ref{ch:s2:capital-structure-arbitrage} --- Capital-Structure Arbitrage}

\begin{solution}{exo:s2:capital-structure-arbitrage:1}
$0.6 \times (-\ln 0.954)/5 = 0.00565$: 56.5 basis points.
\end{solution}

\begin{solution}{exo:s2:capital-structure-arbitrage:2}
$0.00067 \times 10{,}000{,}000 = 6{,}700$ shares, about \$0.34 million at 50.
\end{solution}

\begin{solution}{exo:s2:capital-structure-arbitrage:3}
$4.5 \times 0.0043 = 0.0194$: 193.5 basis points of notional, before the carry received.
\end{solution}

\begin{solution}{exo:s2:capital-structure-arbitrage:4}
Many gaps come from the trader's wrong barrier, not from a mispricing: the firm's true $L$ differs from the assumed 0.5, so the market spread sits
persistently away from the implied one and the gap never closes.
\end{solution}

\begin{solution}{exo:s2:capital-structure-arbitrage:5}
The debt doubles, so the true spread jumps and the protection seller loses on the CDS; the share jumps on the takeover premium, so the short hedge
loses too. The model assumed equity and credit move in opposite directions; a buyout moves them the same way.
\end{solution}

\begin{solution}{exo:s2:capital-structure-arbitrage:6}
Spreads and share prices are loosely correlated day to day, so any one trade's P\&L is noisy and can be large; across hundreds of trades the noise
averages out and the mispricing's reversion remains.
\end{solution}

\begin{solution}{exo:s2:capital-structure-arbitrage:7}
The worst trade improves from $-347$ to $-179$ basis points and the Sharpe ratio rises from 1.40 to 1.66; model error in the barrier remains, and
44\% of trades still lose.
\end{solution}

\begin{solution}{exo:s2:capital-structure-arbitrage:8}
A gap that persists for six months is more likely a wrong input (the barrier, the debt, the vol) or a change in the firm than a mispricing awaiting
correction; in the chapter's simulation the non-converging gaps are the model's errors.
\end{solution}

\begin{solution}{pb:s2:capital-structure-arbitrage:1}
\begin{enumerate}
\item Assets $E + LD$, asset vol $\sigma_E E/V$, default when assets touch $LD$; inputs: share price, equity vol, debt per share, $L$ and recovery.
\item The CDS spread the model gives from the equity: 55.8 basis points with debt of 40 and 115.4 with 80.
\item The model's asset vol falls with the share price.
\item CreditGrades, 135,759 daily spreads on 261 companies; large single-trade losses, a portfolio Sharpe ratio like other fixed-income arbitrage.
\item Shares per unit of CDS notional offsetting the CDS's move with the share price; about 6,700 shares, \$0.34 million.
\item Enter at a 40\% gap in logs; exit below 5\% or after 180 days.
\item Its true barrier and, for three months after a buyout, its debt.
\item The CDS leg: 20.4 of 26.1 basis points a trade.
\item 440 trades, 26.1 basis points each, 48\% losing, worst $-347$, Sharpe ratio 1.40.
\item Their gaps were model error.
\item Strategies needing more intellectual capital earned significant alphas with positively skewed returns.
\item Sharpe ratio 1.66, worst trade $-179$.
\item \textbf{Named result.} The convergence book earns a Sharpe ratio of 1.40 with 48\% of trades losing; protection sellers caught by a leveraged
buyout lost 98 basis points on average, and with buyouts four times as frequent the Sharpe ratio falls to 0.39.
\item Protection sellers lost 98 basis points on average, buyers earned 220.
\item It cannot be observed and differs by firm; the trade takes it as known.
\item Close non-converging gaps, skip firms in play, update debt from filings, cap position size.
\item Debt and vol as known at the time, costs, trades closed at the limit, buyouts in the sample.
\item Debt--equity dislocation after events.
\item It uses Book~6's first-passage survival to turn equity into a spread.
\item That equity and credit will agree again about the company's default risk.
\end{enumerate}
\end{solution}

\begin{solution}{iq:s2:capital-structure-arbitrage:1}
With a structural model: treat equity as a claim on assets, back out asset value and vol from the share price and equity vol, set a default barrier
from the debt, compute survival and convert it into a spread.
\end{solution}

\begin{solution}{iq:s2:capital-structure-arbitrage:2}
Sell protection and short shares at the model's hedge ratio; earn the carry and the convergence. Risks: the model is wrong (barrier, vol, debt), a
buyout or recapitalisation moves equity and credit the same way, borrow is recalled, and the gap widens before it closes.
\end{solution}

\begin{solution}{iq:s2:capital-structure-arbitrage:3}
A leverage shift on the largest positions, a market-wide widening of spreads with equity steady, a change in equity vol, and a borrow recall; with
the book's gross notional and the loss on trades that never converge.
\end{solution}

\begin{solution}{iq:s2:capital-structure-arbitrage:4}
A mispricing reverts; a model error persists or tracks an input. Check whether the gap moves with the firm's own news or filings, test other
barriers, and look at bonds and options of the same firm.
\end{solution}

\begin{solution}{iq:s2:capital-structure-arbitrage:5}
CDS quotes, share prices, equity vol and borrow daily; debt per share from quarterly filings (late); corporate events as announced. The debt is the
input most likely to be stale.
\end{solution}

\begin{solution}{iq:s2:capital-structure-arbitrage:6}
For $X_t = \ln V_t$ with drift $\nu = -\sigma^2/2$ and $x = \ln(V_0/B)$, the reflection principle for Brownian motion with drift gives
$P(\min_{u \le t} X_u > \ln B) = N\big((x + \nu t)/(\sigma\sqrt t)\big) - e^{-2\nu x/\sigma^2}N\big((-x + \nu t)/(\sigma\sqrt t)\big)$, and with
$\nu = -\sigma^2/2$ the factor $e^{-2\nu x/\sigma^2}$ is $e^{x} = V_0/B$.
\end{solution}
